%%%%%%%%%%%%%%%%%
% This is an sample CV template created using altacv.cls
% (v1.7.2, 28 August 2024) written by LianTze Lim (liantze@gmail.com). Compiles with pdfLaTeX, XeLaTeX and LuaLaTeX.
%
%% It may be distributed and/or modified under the
%% conditions of the LaTeX Project Public License, either version 1.3
%% of this license or (at your option) any later version.
%% The latest version of this license is in
%%    http://www.latex-project.org/lppl.txt
%% and version 1.3 or later is part of all distributions of LaTeX
%% version 2003/12/01 or later.
%%%%%%%%%%%%%%%%

%% Use the "normalphoto" option if you want a normal photo instead of cropped to a circle
% \documentclass[10pt,a4paper,normalphoto]{altacv}

\documentclass[10pt,a4paper,ragged2e,withhyper]{altacv}
%% AltaCV uses the fontawesome5 and simpleicons packages.
%% See http://texdoc.net/pkg/fontawesome5 and http://texdoc.net/pkg/simpleicons for full list of symbols.

% Change the page layout if you need to
\geometry{left=1.25cm,right=1.25cm,top=0.5cm,bottom=1.5cm,columnsep=1.2cm}

% The paracol package lets you typeset columns of text in parallel
\usepackage{paracol}
\usepackage[english, russian]{babel}

% Change the font if you want to, depending on whether
% you're using pdflatex or xelatex/lualatex
% WHEN COMPILING WITH XELATEX PLEASE USE
% xelatex -shell-escape -output-driver="xdvipdfmx -z 0" sample.tex
\iftutex 
  % If using xelatex or lualatex:
  \setmainfont{Roboto Slab}
  \setsansfont{Lato}
  \renewcommand{\familydefault}{\sfdefault}
\else
  % If using pdflatex:
  \usepackage[rm]{roboto}
  \usepackage[defaultsans]{lato}
  % \usepackage{sourcesanspro}
  \renewcommand{\familydefault}{\sfdefault}
\fi

% Change the colours if you want to
\definecolor{SlateGrey}{HTML}{2E2E2E}
\definecolor{LightGrey}{HTML}{666666}
\definecolor{DarkPastelRed}{HTML}{450808}
\definecolor{PastelRed}{HTML}{8F0D0D}
\definecolor{GoldenEarth}{HTML}{E7D192}
\colorlet{name}{black}
\colorlet{tagline}{PastelRed}
\colorlet{heading}{DarkPastelRed}
\colorlet{headingrule}{GoldenEarth}
\colorlet{subheading}{PastelRed}
\colorlet{accent}{PastelRed}
\colorlet{emphasis}{SlateGrey}
\colorlet{body}{LightGrey}

% Change some fonts, if necessary
\renewcommand{\namefont}{\Huge\rmfamily\bfseries}
\renewcommand{\personalinfofont}{\footnotesize}
\renewcommand{\cvsectionfont}{\LARGE\rmfamily\bfseries}
\renewcommand{\cvsubsectionfont}{\large\bfseries}


% Change the bullets for itemize and rating marker
% for \cvskill if you want to
\renewcommand{\cvItemMarker}{{\small\textbullet}}
\renewcommand{\cvRatingMarker}{\faCircle}
% ...and the markers for the date/location for \cvevent
% \renewcommand{\cvDateMarker}{\faCalendar*[regular]}
% \renewcommand{\cvLocationMarker}{\faMapMarker*}


% If your CV/résumé is in a language other than English,
% then you probably want to change these so that when you
% copy-paste from the PDF or run pdftotext, the location
% and date marker icons for \cvevent will paste as correct
% translations. For example Spanish:
% \renewcommand{\locationname}{Ubicación}
% \renewcommand{\datename}{Fecha}


%% Use (and optionally edit if necessary) this .tex if you
%% want to use an author-year reference style like APA(6)
%% for your publication list
% \input{pubs-authoryear.tex}

%% Use (and optionally edit if necessary) this .tex if you
%% want an originally numerical reference style like IEEE
%% for your publication list
\input{pubs-num.tex}

%% sample.bib contains your publications
\addbibresource{sample.bib}
% \usepackage{academicons}\let\faOrcid\aiOrcid
\begin{document}
\name{Рауль Алимбеков}
\tagline{Embedded Developer}
%% You can add multiple photos on the left or right
\photoR{2.8cm}{photo}
% \photoL{2.5cm}{Yacht_High,Suitcase_High}

\personalinfo{%
  % Not all of these are required!
  \email{raulalimbekov@gmail.com}
  \phone{+7-987-626-7271}
  \location{Казань, Россия}
  \github{RaulIQ}
  \printinfo{\faTelegram}{@Raul\_no\_Rules}[https://t.me/Raul_no_Rules]

  %% You can add your own arbitrary detail with
  %% \printinfo{symbol}{detail}[optional hyperlink prefix]
  % \printinfo{\faPaw}{Hey ho!}[https://example.com/]

  %% Or you can declare your own field with
  %% \NewInfoFiled{fieldname}{symbol}[optional hyperlink prefix] and use it:
  % \NewInfoField{gitlab}{\faGitlab}[https://gitlab.com/]
  % \gitlab{your_id}
  %%
  %% For services and platforms like Mastodon where there isn't a
  %% straightforward relation between the user ID/nickname and the hyperlink,
  %% you can use \printinfo directly e.g.
  % \printinfo{\faMastodon}{@username@instace}[https://instance.url/@username]
  %% But if you absolutely want to create new dedicated info fields for
  %% such platforms, then use \NewInfoField* with a star:
  % \NewInfoField*{mastodon}{\faMastodon}
  %% then you can use \mastodon, with TWO arguments where the 2nd argument is
  %% the full hyperlink.
  % \mastodon{@username@instance}{https://instance.url/@username}
}

\makecvheader
%% Depending on your tastes, you may want to make fonts of itemize environments slightly smaller
% \AtBeginEnvironment{itemize}{\small}

%% Set the left/right column width ratio to 6:4.
\columnratio{0.65}

% Start a 2-column paracol. Both the left and right columns will automatically
% break across pages if things get too long.
\begin{paracol}{2}


\cvsection{Опыт Работы}

\cvevent{Программист встраиваемых систем}{ЮВС Авиа (концерн "Калашников")}{Сентябрь 2025 — настоящее время}{}

\begin{itemize}
  \item Реализовал \textbf{ППРЧ} на Wi-Fi адаптерах.
  \item Реализовал \textbf{Bidirectional DSHOT} для считывания оборотов моторов.
  
  Проект: \github{RaulIQ/bit-banging-bidir-dshot}
  \item Реализовал мост с аппаратным шифрованием MAVLink (AES-128).
\end{itemize}

\cvevent{Лаборант Центра БПЛА}{МТУСИ}{Март 2024 — Сентябрь 2025}{}
\begin{itemize}
  \item Реализовал протокол \textbf{DSHOT} для управления моторами.
  
  Пример кода: \github{RaulIQ/waveform_example}
  \item Внедрил режимы полёта Acro и Stab на базе PID-регуляторов, цифровых фильтров и кватернионов.
  \item Реализовал запись телеметрии ("черный ящик") на внешнюю SPI-flash с собственным алгоритмом сериализации данных.
  Сделал форк архива spi-memory и внедрил асинхронные API для чтения и записи.
  
  \github{RaulIQ/spi-memory-async}
\end{itemize}

\cvsection{Open Source}
\cvevent{Pull Requests в Embassy}{Добавил новые функциональности в библиотеку с \textbf{9} тысячами звёзд.}{}{}
\begin{itemize}
  \item \github{embassy-rs/embassy/pull/4232}
  \item \github{embassy-rs/embassy/pull/4726}
\end{itemize}

\cvevent{Контрибьютор open-source библиотек}{}{}{}
\begin{itemize}
  \item \github{wboayue/fusion-ahrs}
  \item \github{sulami/dshot-frame}
  \item \github{jettify/uf-ahrs}
\end{itemize}

\cvsection{Проекты}

% может быть обьеденить это все в организацию mik32-rs

\cvevent{mik32-rs}{Набор инструментов для разработки на Rust под микроконтроллеры серии Mik32.}{}{}

\begin{itemize}
\item \github{mik32-rs/mik32-rt}

Стартовый код и минимальный runtime для \textbf{Mik32 Amur}: даёт возможность писать и собирать прошивки на \textbf{Rust} под этот микроконтроллер.

\item \github{mik32-rs/mik32v2-pac}

API периферийного доступа для микроконтроллера mik32v2

\item \github{mik32-rs/mik32-hal}

Предоставляет аппаратную абстракцию, опираясь на трейты \newline \smallskip \textbf{embedded‑hal} и следуя канонам разработки Rust HAL.

\end{itemize}

\medskip

% use ONLY \newpage if you want to force a page break for
% ONLY the current column

%% Switch to the right column. This will now automatically move to the second
%% page if the content is too long.
\switchcolumn

\cvsection{Образование}

\cvevent{Высшее образование}{МТУСИ, Бакалавриат}{Сентябрь 2022 -- Июнь 2026}{}
Четвёртый курс направления: "Искусственный интеллект и машинное обучение"

\smallskip
\textbf{Фокус обучения:} алгоритмы, программная инженерия, ML/AI, DevOps и разработка интеллектуальных систем (CV/SLAM, обработка сигналов).


\cvsection{Навыки}
\textbf{Языки:} \cvtag{Rust} \cvtag{C/C++} \cvtag{Python} \\
\textbf{Встраиваемые платформы:} \cvtag{STM32} \cvtag{ESP32} \cvtag{Mik32} \cvtag{RPi}\\
\textbf{Cреды выполнения:} \cvtag{Bare-metal} \cvtag{FreeRTOS} \cvtag{Embassy (async)} \newline \cvtag{Embedded Linux} \\
\textbf{Коммуникации:} \cvtag{UART/SPI/I2C} \cvtag{SBUS/DSHOT} \cvtag{MAVLink} \\
\textbf{Отладка и инструменты:} \cvtag{GDB} \cvtag{OpenOCD} \cvtag{JTAG/SWD} \cvtag{осциллограф} \cvtag{лог. анализатор} \\



\cvsection{Профессиональная \newline \smallskip активность}

\begin{itemize}
    \item Участвовал в студкемпе "Алиса и Умные устройства Яндекса"
    \item Получал стипендию Selectel Career Wave:
    \href{https://files.selectel.ru/docs/ru/careerwave_scholarship_finalists.pdf}{Список финалистов}
    
    \item \href{https://t.me/it_bash_forum/24}{Выступление на форуме IT-Bash}
\end{itemize}

\cvsection{Интересы и \newline \smallskip увлечения}
\begin{itemize}
  \item Увлекаюсь единоборствами: занимаюсь грэпплингом и бразильским джиу-джитсу в клубе \href{https://combatgym.ru/}{Combat Gym}.
\end{itemize}

\end{paracol}



\end{document}
